% General lifting theorem; external inputs are stated, not reproved.
\section{Lifting a planar cover into three dimensions}\label{sec:lift}
The key construction applies to every convex planar universal cover, and
separates the geometric argument from the choice of a particular planar
cover. Planar congruences in the hypothesis may include reflections; the
resulting three-dimensional cover works using proper rotations alone.

\paragraph{External support-plane theorem.}
We use Wienholtz's alternating-support result in the form proved by
Ghomi~\cite[Corollary~3.2]{Ghomi2018}.
For a continuous space curve whose affine span is three-dimensional, there
are two parallel supporting planes and four parameter values
$t_0<t_1<t_2<t_3$ at which the curve meets the two planes alternately.
The curve lies in the slab between those planes. For a curve contained in a
plane, we will instead use that containing plane as a slab of width zero.
This is a nontrivial external theorem; the argument below is a deduction
from it. The planes need not realize the curve's minimum width.

\begin{theorem}[A planar cover produces a spatial cover]\label{thm:lift}
Let $C\subset\mathbb R^2$ be a compact convex set with nonempty interior,
containing the origin, which contains a congruent copy of every planar
rectifiable curve of length at most one. Set
\[
 f(z)=\sqrt{1-36z^2},\qquad -\frac16\le z\le\frac16,
\]
and define
\[
 K_C=\left\{(p,z)\in\mathbb R^2\times\mathbb R:
 |z|\le\frac16,\quad p\in f(z)C\right\}.
\]
Then $K_C$ is a convex body which contains a translated and properly rotated
copy of every rectifiable space curve of length at most one. Moreover,
\[
 \boxed{|K_C|=\frac29\operatorname{area}(C).}
\]
\end{theorem}
The assumption $0\in C$ entails no restriction on the area of a planar
cover: any chosen point of the cover may be translated to the origin.

\paragraph{Convexity and compactness.}
The function $f$ is continuous and concave on $[-1/6,1/6]$. Compactness of
$K_C$ follows from its representation as the continuous image of
$C\times[-1/6,1/6]$ under $(c,z)\mapsto(f(z)c,z)$. To check convexity,
take $(p_i,z_i)\in K_C$, write $p_i=f(z_i)c_i$ with $c_i\in C$, and let
$0\le\lambda\le1$. Put
\[
 z=(1-\lambda)z_1+\lambda z_2,\qquad
 a=(1-\lambda)f(z_1)+\lambda f(z_2).
\]
Convexity of $C$ gives
$(1-\lambda)p_1+\lambda p_2\in aC$; when $a=0$ this simply says that the
point is zero. Concavity gives $a\le f(z)$, and $0\in C$ implies
$aC\subseteq f(z)C$. Thus the convex combination belongs to $K_C$.
The product $(C/2)\times[-\sqrt3/12,\sqrt3/12]$ is contained in $K_C$,
so $K_C$ has nonempty interior even when the origin is a boundary point
of $C$.

\paragraph{Length of the two projections.}
Fix a curve $\gamma$ of length $L\le1$. A constant curve is immediate, so
assume $L>0$. Apply the support-plane theorem if the curve is nonplanar;
otherwise take a plane containing the curve. By a proper rigid motion,
put the resulting supporting planes at $z=-h/2$ and $z=h/2$, where $h\ge0$.
Write the transformed curve as $\gamma(s)=(p(s),z(s))$ and parametrize it
by arclength. Let $\ell$ be the length of the planar projection $p$, and
let $v$ be the total variation of $z$. The three alternating traversals
between the supporting planes give
\[
 v\ge3h.
\]
This inequality also holds in the planar case $h=0$.
Because an arclength parametrization is Lipschitz, its derivative exists
almost everywhere and the projection lengths satisfy
\[
 \ell=\int_0^L |p'(s)|\,ds,\qquad
 v=\int_0^L |z'(s)|\,ds.
\]
The triangle inequality for the integral of the two-dimensional vector
$(|p'(s)|,|z'(s)|)$ yields
\[
 \sqrt{\ell^2+v^2}
 \le \int_0^L\sqrt{|p'(s)|^2+|z'(s)|^2}\,ds=L\le1.
\]
Consequently
\[
 0\le h\le\frac13,
 \qquad \ell\le\sqrt{1-9h^2}.
\]
This is the essential link between the thickness of the chosen slab and
the amount of length left in its orthogonal projection. The quantity $h$
is the width of this chosen slab, not necessarily the minimum width over
all directions.

\paragraph{One rigid placement fits every height.}
If $\ell>0$, the planar universality of $C$, applied after scaling the
projection by $1/\ell$, supplies $A\in O(2)$ and $b\in\mathbb R^2$ such that
\[
 A p(s)+b\in\ell C\quad\hbox{for every }s.
\]
If $\ell=0$, the projection is a point; a translation sends it to the
origin, giving the same conclusion with $\ell C=\{0\}$.
For all points of the curve,
\[
 |z(s)|\le h/2\le1/6,
 \qquad
 f(z(s))\ge\sqrt{1-9h^2}\ge\ell.
\]
Since $0\in C$, we have $\ell C\subseteq f(z(s))C$. Hence this single
planar placement fits the entire three-dimensional curve into $K_C$ at
all heights simultaneously.

For completeness, even if the planar placement uses a reflection, the
three-dimensional placement can be chosen proper: put
$\varepsilon=\det A\in\{-1,1\}$ and use
\[
 (p,z)\longmapsto(Ap+b,\varepsilon z).
\]
Its linear part has determinant $(\det A)\varepsilon=1$, while replacing
$z$ by $-z$ preserves every slice of $K_C$. Composing with the initial
proper rigid motion proves the assertion without permitting reflections
of the physical space curve.

\paragraph{Volume.}
The section at height $z$ has area $f(z)^2\operatorname{area}(C)$. Thus
Cavalieri's principle gives
\[
 |K_C|=\operatorname{area}(C)
 \int_{-1/6}^{1/6}(1-36z^2)\,dz
 =\frac29\operatorname{area}(C).
\]

\subsection{An explicit ellipsoidal sector of volume \texorpdfstring{$\pi/54$}{pi/54}}
The planar theorem of Panraksa and Wichiramala~\cite{PW2021}
states that a circular sector of unit radius
and angle $\pi/6$ contains a congruent copy of every unit-length planar arc.
The statement for shorter arcs follows by extending the path to unit
length. Apply the lifting theorem to the sector
\[
 S=\{(\rho\cos\theta,\rho\sin\theta):
 0\le\rho\le1,\ |\theta|\le\pi/12\},
 \qquad \operatorname{area}(S)=\frac{\pi}{12}.
\]
Since $\tan(\pi/12)=2-\sqrt3$, its lift is the explicit convex body
\[
 \boxed{K_S=\left\{(x,y,z):
 x^2+y^2+36z^2\le1,\quad
 x\ge0,\quad |y|\le(2-\sqrt3)x\right\}.}
\]
It is the intersection of an ellipsoid with a $30$-degree dihedral wedge,
and
\[
 \boxed{|K_S|=\frac29\frac{\pi}{12}
 =\frac{\pi}{54}=0.05817764173314\ldots.}
\]
This is a rigorously justified feasible volume; the lifting argument does
not establish that it is the minimum. The more general theorem also shows
how any improved convex planar cover immediately gives a corresponding
three-dimensional upper bound.

\begin{figure}[htbp]
\centering
\includegraphics[width=\textwidth]{cover.pdf}
\caption{Left: the explicit universal convex body, an ellipsoid intersected with a $30^\circ$ dihedral wedge. Right: for any chosen alternating-support slab, one prism of height $h$ and planar scale $\ell$ fits inside the body. The figure illustrates the proof; containment is proved by the exact inequality $\ell^2+9h^2\le1$.}
\end{figure}
